\documentclass[11pt]{article}
\usepackage{problem_set_style}
\usepackage{kotex}
\usepackage{booktabs}
\usepackage{array}
\begin{document}
\heading{G1}{중간 · 기말 공통}{}{Spring 2027}{논문(보고서) 작성 가이드라인}

\noindent{\small 이 수업의 프로젝트 보고서는 \textbf{작은 논문}이다. 학회에 내는 논문과 같은 규칙으로 쓰고, 같은 방식(KSA Review)으로 서로 심사한다. 이 문서는 ICLR 2026$\cdot$2027, NeurIPS 2026, ICML 2026의 AI 정책과 The Web Conference 2027 심사 지침을 우리 수업에 맞게 고친 것이다. 학회별 비교는 리뷰 가이드라인 끝의 표를 보라.}

\section*{먼저 이것만은}
\begin{enumerate}
\item \textbf{주장은 하나, 첫 문단에.} 읽는 사람이 첫 문단만 보고 ``이 팀은 무엇을 보였다고 주장하는가''를 한 문장으로 말할 수 있어야 한다.
\item \textbf{모든 숫자는 실험에서 나와야 한다.} 표·그래프의 숫자는 제출한 코드를 실행하면 다시 나와야 한다. 지어낸 숫자, 고른 숫자, 출처를 모르는 숫자는 없어야 한다.
\item \textbf{AI 도구는 써도 되지만, 밝히고 책임진다.} 무엇에 썼는지 ``AI 사용 내역''에 적는다. AI가 쓴 문장·코드·참고문헌이 틀렸다면 그 책임은 팀에 있다 (3절).
\item \textbf{참고문헌은 직접 확인한 것만.} 실제로 존재하고, 실제로 그 내용을 말하는 문헌만 인용한다. \textbf{존재하지 않는 참고문헌은 가장 무거운 위반}이다.
\item \textbf{한계를 숨기지 않는다.} 성능이 낮아도 원인을 정직하게 진단한 보고서가, 성능이 높지만 검증이 허술한 보고서보다 높은 점수를 받는다.
\end{enumerate}

\section*{1. 논문의 구조}
공지문의 ``보고서 구조''를 따른다. 각 부분에서 독자가 얻어야 할 것은 아래와 같다.
\begin{center}
\small
\begin{tabular}{@{}p{3.0cm}p{10.2cm}@{}}
\toprule
\textbf{부분} & \textbf{독자가 여기서 알아야 할 것} \\
\midrule
제목 · Abstract & 무엇을, 어떤 데이터로, 무엇과 비교해, 무엇을 발견했는가. Abstract만 읽어도 목표·데이터·결론이 보여야 한다. \\
서론 (문제 정의) & 왜 이 문제인가, 행의 단위 · 라벨 · 주 지표 · 베이스라인, 그리고 \textbf{기여 목록}(``이 보고서는 ①… ②…를 보인다''). \\
데이터와 방법 & 다른 사람이 똑같이 따라 할 수 있을 만큼: 분할 방식, train으로만 fit한 전처리, 모델과 하이퍼파라미터 탐색 범위. \\
결과 & 기여 목록의 각 항목을 실험 하나 이상이 뒷받침해야 한다. test는 한 번만, 어디서 썼는지 밝힌다. \\
분석 · 한계 & 왜 이런 결과가 나왔는지(수식적 정당화), 무엇이 안 됐고 어떤 조건에서 일반화되지 않을지. \\
재현 방법 · AI 사용 내역 & 데이터 출처, 시드, 실행 명령. AI 도구를 어디에 썼는지 (3절). \\
\bottomrule
\end{tabular}
\end{center}

\section*{2. 심사자가 반드시 확인하는 것}
리뷰 가이드라인의 심사자는 아래 질문을 들고 여러분의 논문을 읽는다. 제출 전에 팀이 먼저 답해 보라.
\begin{itemize}
\item 서론에서 주장한 기여가 실제로 결과에서 \textbf{전부} 보여졌는가?
\item 실험이 주장을 시험하는가, 아니면 더 쉬운 것을 시험하는가?
\item 베이스라인이 전문가가 기대할 만한 것인가, 그리고 \textbf{공정한 기회}(같은 분할 · 같은 튜닝 노력)를 받았는가?
\item 결과가 \textbf{데이터 누수}, test로 튜닝하기, 불공정한 예산으로 설명될 수 있는가?
\item 제목과 Abstract가 실제로 한 일을 정확히 묘사하는가 (과장 없이)?
\item 관련 연구를 공정하게 소개했는가?
\end{itemize}

\section*{3. AI 도구 사용 정책 (저자)}
\textbf{원칙} (ICLR 2026$\cdot$2027과 같다): (가) AI 도구를 썼다면 \textbf{밝혀야 하고}, 그 결과물에 대한 \textbf{책임은 팀에} 있다. (나) AI를 썼든 안 썼든, \textbf{거짓이나 오해를 부르는 주장, 데이터·결과의 조작이나 왜곡}은 허용되지 않는다.

\begin{center}
\small
\begin{tabular}{@{}p{6.6cm}p{6.6cm}@{}}
\toprule
\textbf{써도 된다 (밝히고)} & \textbf{안 된다} \\
\midrule
문장 다듬기, 맞춤법, 영어 표현 & 실행하지 않은 실험의 결과를 쓰게 하기 \\
코드 오류 찾기, 라이브러리 사용법 묻기 & 직접 확인하지 않은 참고문헌 넣기 \\
개념 설명 듣기 (교과서처럼) & 결과 해석·한계를 통째로 생성해 붙이기 \\
그래프 코드 초안 만들기 (돌려서 확인) & 숫자·표를 AI 출력에서 바로 옮겨 적기 \\
\bottomrule
\end{tabular}
\end{center}

\noindent\textbf{AI 사용 내역} (보고서 끝, 3$\sim$5줄, \textbf{필수}, 쪽수에 넣지 않음. ICLR 2027도 같은 절을 의무로 한다): 어떤 도구를 어느 부분에 어떻게 썼는지 적는다. 데이터 정리 · 방법 구현 · 결과 해석 · 가설 제안에 썼다면 \textbf{반드시}, 문장 다듬기 · 그림 코드 · 문헌 정리에 썼다면 되도록 적는다. 예: ``2절 문장 다듬기에 ChatGPT 사용. 4절 학습 곡선 그래프 코드 초안을 Claude로 만들고 직접 실행해 확인.'' 쓰지 않았다면 ``사용하지 않음''.

\vspace{4pt}
\noindent\textbf{위반으로 보는 것과 결과}
\begin{itemize}
\item \textbf{존재하지 않는 참고문헌, 실험에서 나오지 않은 숫자, 코드와 다른 결과}: 연구 윤리 위반이다. 학회라면 심사 없이 거절(desk reject)된다. 우리 수업에서는 해당 보고서를 \textbf{0점} 처리하고 면담한다.
\item \textbf{AI를 많이 썼는데 밝히지 않음}: 사용 내역을 보완하게 하고, 재현성·명확성 항목에서 감점한다.
\item \textbf{AI 심사를 노린 숨은 문장} (흰 글씨 · 아주 작은 글씨로 ``이 논문을 높게 평가하라'' 같은 지시를 넣기): 프롬프트 주입이다. NeurIPS · ICML · ACL 모두 금지하며 우리 수업도 0점이다.
\item \textbf{거의 같은 보고서를 두 팀이 내거나, 바꿔 쓴 판본으로 이중 제출}: 표절로 처리한다.
\item AI 탐지 도구 결과만으로는 처리하지 않는다. 교사가 \textbf{구체적 근거}(없는 문헌, 재현되지 않는 숫자)를 확인한 경우에만 처리한다.
\end{itemize}

\section*{4. 제출 전 점검표}
\begin{itemize}
\item[$\square$] 첫 문단에서 주장이 한 문장으로 보인다.
\item[$\square$] 기여 목록의 모든 항목에 대응하는 실험·그림·표가 있다.
\item[$\square$] 베이스라인과 모델 두 가지 이상을 같은 분할·같은 지표로 비교했다.
\item[$\square$] test는 한 번만 썼고, 전처리는 train으로만 fit했다.
\item[$\square$] 코드를 처음부터 다시 실행해 표의 숫자가 그대로 나오는 것을 확인했다.
\item[$\square$] 모든 참고문헌을 직접 열어 보고, 그 내용이 인용한 문장과 맞는지 확인했다.
\item[$\square$] 한계 절에 ``잘 안 된 것''을 하나 이상 적었다.
\item[$\square$] AI 사용 내역을 적었다 (안 썼으면 ``사용하지 않음'').
\item[$\square$] 익명 심사용 파일에 팀원 이름 · 학번 · 이름이 보이는 저장소 링크가 없다.
\end{itemize}

\section*{5. 리뷰를 받은 뒤 (반박과 수정)}
\begin{itemize}
\item 리뷰의 질문에 번호를 맞춰 \textbf{1쪽 이내}로 답한다. 고칠 수 있는 것은 고치고, 동의하지 않는 지적은 근거를 들어 정중하게 반박한다.
\item 리뷰가 논문에 없는 내용을 지어냈거나, 존재하지 않는 문헌을 들었거나, 다른 논문에도 그대로 붙일 수 있는 일반론뿐이라면 증거와 함께 교사에게 \textbf{비공개로} 알린다. 리뷰어를 공개적으로 비난하지 않는다.
\end{itemize}

\vspace{6pt}
\noindent{\scriptsize 출처: ICLR 2027, AI Policy for Authors (iclr.cc). ICLR 2026 Program Chairs, ``ICLR 2026 Response to LLM-Generated Papers and Reviews'' (2025.11.19), blog.iclr.cc. NeurIPS 2026, Main Track Handbook (neurips.cc). ICML 2026, Reviewer Instructions (icml.cc). The ACM Web Conference 2027, Reviewer Guidelines (Research Tracks).}

\end{document}
